\chapter{Sputum}

\section*{Definition}
Sputum is mucus and other matter expectorated from the lower respiratory tract (tracheobronchial tree), produced in response to inflammation, infection, or irritant exposure. Normal airway secretions are cleared subconsciously; visible sputum production indicates excessive secretion or impaired mucociliary clearance.

\section*{Pathophysiology}
Airway irritation and inflammation trigger goblet cell hyperplasia and submucosal gland hypertrophy, increasing mucus production. Inflammatory mediators (histamine, leukotrienes, prostaglandins) stimulate mucin gene expression and secretion. Impaired ciliary function from smoke, infection, or genetic conditions (cystic fibrosis, primary ciliary dyskinesia) prevents mucus clearance, leading to accumulation and expectation. Neutrophil recruitment and degranulation contribute to purulence and color changes.

\section*{Etiology}
\begin{itemize}
  \item Acute infections: Acute bronchitis, Pneumonia (bacterial, viral, fungal), Tracheitis, Pertussis
  \item Chronic airway diseases: Chronic obstructive pulmonary disease (COPD), Asthma, Bronchiectasis, Cystic fibrosis
  \item Irritant exposure: Cigarette smoke, Occupational dusts (coal, silica, asbestos), Air pollution, Aspiration (gastric contents, foreign body)
  \item Pulmonary infections: Tuberculosis, Nontuberculous mycobacteria, Fungal infections (aspergillosis, histoplasmosis), Parasitic infections
  \item Neoplastic causes: Lung cancer, Bronchial adenoma, Lymphangitic carcinomatosis
  \item Other causes: Pulmonary edema (heart failure), Gastroesophageal reflux disease, Post-nasal drip, Bronchopleural fistula
\end{itemize}

\section*{Indications for Evaluation}
Seek care if:
\begin{itemize}
  \item Sputum changes color (green, rust, bloody) or increases in volume
  \item Hemoptysis (blood-streaked sputum or frank blood)
  \item Sputum production with fever, chills, pleuritic chest pain, or dyspnea
  \item Chronic productive cough for $>$8 weeks
  \item Immunocompromised patient with new sputum production
\end{itemize}

\section*{Related Symptoms}
\begin{itemize}
  \item Cough
  \item Dyspnea
  \item Wheezing
  \item Chest pain
  \item Fever
  \item Hemoptysis
  \item Nasal congestion
  \item Post-nasal drip
\end{itemize}
\textbf{Note:} Sputum Gram stain and culture guide antibiotic therapy in pneumonia; acid-fast staining and culture are indicated when tuberculosis is suspected.

\section*{Self-Care / Home Management}
\begin{itemize}
  \item Stay well hydrated to thin mucus and facilitate expectation.
  \item Use a humidifier or steam inhalation to moisten airways and soothe irritation.
  \item Practice controlled coughing (huff coughing) to clear secretions effectively.
  \item Avoid smoking and exposure to airborne irritants during recovery.
  \item Over-the-counter expectorants (guaifenesin) may assist mucus clearance.
\end{itemize}

\section*{Diagnostic Approach}
\begin{itemize}
  \item Assess sputum color, volume, consistency, and presence of blood.
  \item Sputum Gram stain and culture with sensitivity for suspected pneumonia or bronchitis.
  \item Acid-fast bacillus smear, culture, and PCR for tuberculosis.
  \item Cytology for suspected malignancy in high-risk patients (smokers, hemoptysis).
  \item Chest X-ray or CT for suspected pneumonia, bronchiectasis, or lung mass.
  \item Spirometry and pulmonary function testing for chronic productive cough.
\end{itemize}

\section*{Treatment / Management Overview}
\begin{itemize}
  \item Acute infections: Antibiotics for bacterial pneumonia/bronchitis based on culture results; supportive care for viral infections.
  \item COPD/bronchiectasis: Bronchodilators, inhaled corticosteroids, pulmonary rehabilitation, airway clearance techniques.
  \item Cystic fibrosis: Mucolytics (dornase alfa), hypertonic saline, CFTR modulators, airway clearance.
  \item Smoking cessation counseling and support for all active smokers.
\end{itemize}

\clearpage
